% REPLACEMENT TEXT ONLY: replace Sections II.D--II.G of PRL_SSG_Type_sup.
% Keep Sections II.A--II.C and all of Section III in the formal SM.
% Requires amsmath, amssymb, bm and hyperref.
% Original equation numbers (9)--(38) are retained. Added equations have
% letter suffixes. All new labels have the prefix algmerge:.
% The counters below allow this file to be inserted after Section II.C.
\setcounter{section}{2}
\setcounter{subsection}{3}
\setcounter{equation}{8}
\begingroup
% Keep hyperlink destinations distinct even though explicit tags preserve
% the formal SM's equation numbers. Restore the normal anchor style below.
\providecommand{\theHequation}{\arabic{equation}}
\renewcommand{\theHequation}{algmerge.\arabic{equation}}

\subsection{Type 0}
\label{algmerge:type0}

By definition, Type 0 SSGs have $\rho_{\mathrm T}(\mathrm T)=\{E\}$.
The normalizer is $N_{SO(3)}(\{E\})=SO(3)$, and $\rho(G)$ is the
image of the finite point group $P$, hence a finite subgroup of $SO(3)$.
The homomorphism is therefore determined by its point-group part:
choosing $\rho(t s(p))=\rho_P(p)$ gives
\stepcounter{equation}
\begin{equation}
 \left.\operatorname{Hom}(G,SO(3))\right|_{\rho_{\mathrm T}=1}
 \simeq\operatorname{Hom}(P,SO(3)).
 \tag{9}\label{algmerge:type0hom}
\end{equation}
Combining this result with the $\mathbb Z_2$ grading gives
\stepcounter{equation}
\begin{equation}
 \left.\operatorname{Hom}(G,O(3))\right|_{\rho_{\mathrm T}=1}
 \simeq\operatorname{Hom}(P,SO(3))\times
 \operatorname{Hom}(G,\mathbb Z_2^{\mathcal T}).
 \tag{10}\label{algmerge:type0graded}
\end{equation}
When $\chi_{\mathrm T}$ is trivial, so that $\varrho_{\mathrm T}=1$,
this reduces to
\stepcounter{equation}
\begin{equation}
 \left.\operatorname{Hom}(G,O(3))\right|_{\varrho_{\mathrm T}=1}
 \simeq\operatorname{Hom}(P,SO(3))\times
 \operatorname{Hom}(P,\mathbb Z_2^{\mathcal T})
 \simeq\operatorname{Hom}(P,O(3)),
 \tag{11}\label{algmerge:spinpoint}
\end{equation}
which is the representation problem for spin point groups associated
with $P$.

To enumerate the inequivalent $\rho_P$, we fix a normal subgroup
$H\lhd P$ and set $\ker\rho_P=H$. The fundamental theorem of
homomorphisms gives
\stepcounter{equation}
\begin{equation}
 \rho_P(P)\simeq P/\ker\rho_P=P/H.
 \tag{12}\label{algmerge:quotient}
\end{equation}
The possible rotation images are $C_n$ with $n=1,2,3,4,6$,
$D_n$ with $n=2,3,4,6$, and the tetrahedral and octahedral groups
$T$ and $O$. No icosahedral image is possible, since the order of a
crystallographic point group is not divisible by five.

The normal subgroups can be obtained by starting from $\{e\}$ and
successively adjoining the normal closure of each additional element.
Repeating this construction gives every normal subgroup of the finite
group $P$. For each $H$, the coset multiplication defines the finite
group $Q=P/H$. To test an isomorphism $f:Q\to S_i$, choose generators
$q_1,\ldots,q_r$ of $Q$ and enumerate their possible images in $S_i$.
Each image must have the same order as its corresponding generator,
and the images must generate $S_i$. Propagate an assignment by group
multiplication; if two words representing the same element of $Q$
give different images, discard it. A surviving assignment is accepted
only when it is bijective and satisfies
\stepcounter{equation}
\begin{equation}
 f(q q')=f(q)f(q')\qquad\text{for every }q,q'\in Q.
 \tag{12a}\label{algmerge:isomorphism}
\end{equation}
Thus equality of group orders or element-order distributions is used
only to exclude candidates, not to establish an isomorphism.

For the rotation groups in the above list, every automorphism is
induced by conjugation in their $SO(3)$ normalizer. For $C_n$, the
allowed changes of a generator are inversion and the identity; for
$D_2$, all permutations of the three half-turn axes are allowed. For
$D_3,D_4,D_6$, inversion of the principal rotation and shifts of the
transverse axes are likewise realized by spin rotations. The
tetrahedral automorphisms are realized by $O$, while the octahedral
automorphisms are inner. Consequently, two isomorphisms to the same
standard image give spin-conjugate homomorphisms. It suffices to keep
one isomorphism for each allowed quotient and to pull it back to $P$.
This conclusion uses the particular list of crystallographic rotation
images; it is not a general assertion that homomorphisms with equal
kernels are conjugate.

The enumeration therefore consists of enumerating $H\lhd P$, testing
$P/H$ against these rotation groups, and lifting one accepted
isomorphism to $\rho_P$. Different kernels are retained at this stage.
They can be identified by spatial automorphisms only after $\rho_P$
has been combined with $\chi$, as described in
Sec.~\ref{algmerge:equivalences}.

\subsection{Type 1}
\label{algmerge:type1}

\subsubsection{Translation holonomy}

For Type 1 SSGs, the image $\rho_{\mathrm T}(\mathrm T)\subset
SO(2)_{\bm n}$ is nontrivial and consists of rotations about a common
axis $\bm n$. Choose $\bm n=\hat{\bm z}$ and write
\stepcounter{equation}
\begin{equation}
 \rho_{\mathrm T}(\bm t)=e^{-i\theta(\bm t)J_z}.
 \tag{13}\label{algmerge:axialtranslation}
\end{equation}
The translation homomorphism condition, Eq.~(5), requires
\stepcounter{equation}
\begin{equation}
 \theta(\bm t_1)+\theta(\bm t_2)=\theta(\bm t_1+\bm t_2)
 \pmod{2\pi},\qquad \bm t_1,\bm t_2\in L.
 \tag{14}\label{algmerge:additive}
\end{equation}
Hence $\theta(\bm t)=\bm k\cdot\bm t$ for a wave vector defined
modulo the reciprocal lattice $\widehat L$, and
\stepcounter{equation}
\begin{equation}
 \rho_{\mathrm T}(\bm t)=e^{-i(\bm k\cdot\bm t)J_z}.
 \tag{15}\label{algmerge:wavevector}
\end{equation}
The value $\bm k\equiv0$ is excluded from Type 1 because it gives
Type 0. It will nevertheless be retained as an intermediate parameter
when studying the continuation of a general representation branch.

\subsubsection{Decomposition of \texorpdfstring{$\rho$}{rho} and the covariant condition}

The normalizer of the nontrivial axial translation image is
\stepcounter{equation}
\begin{equation}
 N_{SO(3)}\bigl(\rho_{\mathrm T}(\mathrm T)\bigr)
 =SO(2)_z\rtimes\{E,R_x(\pi)\}\simeq SO(2)\rtimes C_2.
 \tag{16}\label{algmerge:axialnormalizer}
\end{equation}
This holds both for a finite cyclic image, including $C_2$, and for
an infinite image dense in $SO(2)$. The section matrices therefore
take the form
\stepcounter{equation}
\begin{equation}
 \rho(s(p))=e^{-i\theta(s(p))J_z}[R_x(\pi)]^{(1-\eta_p)/2},
 \qquad\eta:P\to\{\pm1\},
 \tag{17}\label{algmerge:axialsection}
\end{equation}
where $\eta_p$ specifies whether the orientation of the spin axis is
preserved or reversed. Substitution into Eq.~(6) gives, in Cartesian
coordinates,
\stepcounter{equation}
\begin{equation}
 e^{-i\eta_p(\bm k\cdot\bm t)J_z}
 =e^{-i(R_p^{-1}\bm k\cdot\bm t)J_z},
 \tag{18}\label{algmerge:covariancematrix}
\end{equation}
or, equivalently,
\stepcounter{equation}
\begin{equation}
 R_p\bm k=\eta_p\bm k\pmod{\widehat L},\qquad p\in P.
 \tag{19}\label{algmerge:covariance}
\end{equation}
At a nonzero TRIM, $\bm k\equiv-\bm k$, so every point-group
character $\eta$ is compatible with an invariant $\bm k$ at the
level of this covariance condition. The section equation may impose
a further obstruction. At a non-TRIM wave vector, an allowed
$\bm k$ fixes $\eta$ uniquely. We therefore enumerate $\eta$ first
and solve for its compatible wave vectors.

\subsubsection{Obstructed 1-cocycle equation of point-group sections}

Substituting Eq.~\eqref{algmerge:axialsection} into Eq.~(7) gives
\stepcounter{equation}
\begin{equation}
 \theta(s(p_1))+\eta_{p_1}\theta(s(p_2))
 =\bm k\cdot\bm t_{p_1,p_2}+\theta(s(p_1p_2))\pmod{2\pi}.
 \tag{20}\label{algmerge:type1cocycle}
\end{equation}
Equivalently, $e^{-i\theta(s(p))}K^{(1-\eta_p)/2}$ defines a
one-dimensional corepresentation. Its factor system is
$\nu_{\bm k}(p_1,p_2)=e^{-i\bm k\cdot\bm t_{p_1,p_2}}$.
A solution exists precisely when this factor system is a coboundary.
When it exists, the solutions modulo phase coboundaries form a torsor
for $H^{1,\eta}(P,U(1))$: choosing one particular solution identifies
the remaining freedom with this cohomology group.

This statement provides a finite method of solving the obstruction.
Put $x_p=\theta(s(p))/(2\pi)$ and
$f_{p_1p_2}=\bm k\cdot\bm t_{p_1,p_2}/(2\pi)$.
The integer matrix with rows indexed by $(p_1,p_2)$ is
\stepcounter{equation}
\begin{equation}
 (A_\eta)_{(p_1,p_2),r}
 =\delta_{r,p_1}+\eta_{p_1}\delta_{r,p_2}-\delta_{r,p_1p_2},
 \qquad A_\eta x=f\pmod1.
 \tag{20a}\label{algmerge:phaselinear}
\end{equation}
A common spin rotation changes
$x_p$ to $x_p+(1-\eta_p)v$. If $\eta$ is nontrivial, choose
$r$ with $\eta_r=-1$ and impose $x_r=0$, or equivalently use
\stepcounter{equation}
\begin{equation}
 x_p\longmapsto x_p-\frac{1-\eta_p}{2}x_r.
 \tag{20b}\label{algmerge:gauge}
\end{equation}
This removes the continuous coboundary before counting solutions.
For trivial $\eta$ there is no such gauge direction.

With the gauge row included when needed, let $U A_\eta V=D$ be
the Smith form. Here $U,V$ are integer unimodular, and $D$ has
positive diagonal entries $d_1,\ldots,d_r$, with all other
entries zero. In the coordinates $x=Vy$,
all solutions are given by
\stepcounter{equation}
\begin{equation}
 y_i=\frac{(Uf)_i+m_i}{d_i}\pmod1,\quad
 m_i=0,\ldots,d_i-1\quad(i\le r),
 \qquad (Uf)_i=0\pmod1\quad(i>r).
 \tag{20c}\label{algmerge:smithfixed}
\end{equation}
The conditions on zero rows test the obstruction. Free column
coordinates, when present, range over $\mathbb R/\mathbb Z$.
At fixed $\bm k$ the gauge-fixed phase solutions are finite. Thus
no separate search through projective corepresentation tables is
required. Integer, rather than rational, row operations are necessary
to retain the torsion branches. Changing the section changes the
factor system by a coboundary and only reparametrizes these solutions.

Let $\mathcal S_{\eta,\bm k}$ denote this set of phase solutions
modulo the rotation gauge. The Type 1 spin-conjugacy classes are
\stepcounter{equation}
\begin{equation}
 \left(\coprod_{\eta\in\operatorname{Hom}(P,\mathbb Z_2)}
 \coprod_{\bm k\in W_\eta}\mathcal S_{\eta,\bm k}\right)/\mathbb Z_2,
 \tag{21}\label{algmerge:axialclasses}
\end{equation}
where
\stepcounter{equation}
\begin{equation}
 W_\eta=\{\bm k\ne0:\ R_p\bm k\equiv\eta_p\bm k\ (p\in P),
 \ [\nu_{\bm k}]=1\}.
 \tag{22}\label{algmerge:allowedwavevectors}
\end{equation}
The remaining $\mathbb Z_2$ is conjugation by $R_x(\pi)$, acting
as $(\bm k,\theta)\mapsto(-\bm k,-\theta)$. Although each
$\mathcal S_{\eta,\bm k}$ is a torsor for the same $H^{1,\eta}$,
its identification with that group need not be globally consistent
as $\bm k$ winds around a reciprocal period.

\subsubsection{Continuous families and incommensurate orders}

The allowed wave vectors may vary continuously. We retain this
variation as a parameter, allowing both commensurate and
incommensurate orders without a bound on their magnetic periods.
For $\bm k=\sum_j q_j\bm b_j$, the order is commensurate if all
$q_j$ are rational. The continuous classification must retain the
section phases together with the wave vector, since a reciprocal
winding can interchange phase branches.

In a primitive basis, write $M_p$ for the integer matrix of $R_p$
and $\bm t_{p_1,p_2}=\sum_j c_j(p_1,p_2)\bm a_j$.
Equations~\eqref{algmerge:covariance} and
\eqref{algmerge:type1cocycle} become the joint integer system
\stepcounter{equation}
\begin{equation}
 (M_p^T-\eta_p I)\bm q=0\pmod1,\qquad
 x_{p_1}+\eta_{p_1}x_{p_2}-x_{p_1p_2}
 -\bm q\cdot\bm c(p_1,p_2)=0\pmod1.
 \tag{22a}\label{algmerge:joint}
\end{equation}
The reciprocal action in these coordinates is $M_p^{-T}$;
Eq.~\eqref{algmerge:joint} is the equivalent condition obtained
directly by pairing with $M_p\bm n$. Adding the gauge condition
and applying the same Smith procedure gives
\stepcounter{equation}
\begin{equation}
 X_\eta=\coprod_a C_a,\qquad
 C_a=z_a+B_a\mathbb T^{d_a},\qquad
 z=(\bm q,x),\quad\mathbb T=\mathbb R/\mathbb Z.
 \tag{22b}\label{algmerge:jointcomponents}
\end{equation}
The offsets $z_a$ are rational and $B_a$ is integral. This joint
solution preserves the winding of the phase branches. Its use for
continuous equivalence, including the treatment of special points,
is specified in Sec.~\ref{algmerge:continuous}.

\subsubsection{Algorithm}

For a given parent group, the Type 1 enumeration proceeds as follows:
\begin{enumerate}
\renewcommand{\labelenumi}{(\roman{enumi})}
\item Enumerate the point-group characters $\eta$, including the
trivial one. The nontrivial characters are obtained from index-two
normal subgroups of $P$.
\item For each $\eta$, solve the joint wave-vector and section-phase
equations~\eqref{algmerge:joint}, with the gauge fixed as in
Eq.~\eqref{algmerge:gauge}, by Smith decomposition.
\item Intersect the resulting tori with $2\bm q=0$ modulo one.
Retain nonzero TRIM solutions separately, account for $\bm q=0$
in Type 0, and retain the remaining part of each positive-dimensional
component as a general family. Isolated non-TRIM solutions remain
discrete entries.
\item Construct $\rho(t s(p))=\rho_{\mathrm T}(t)\rho(s(p))$ and
identify spin-conjugate solutions, including the simultaneous sign
reversal of wave vector and section phases.
\item Combine with $\chi$ and impose the spatial equivalences of
Sec.~\ref{algmerge:equivalences} on the complete representation.
\end{enumerate}

\subsection{Type 2}
\label{algmerge:type2}

\subsubsection{Translation holonomy}

For Type 2, the translation image is the Klein four group. Choose
\stepcounter{equation}
\begin{equation}
 V_4=\{E,R_x(\pi),R_y(\pi),R_z(\pi)\}\simeq D_2.
 \tag{23}\label{algmerge:klein}
\end{equation}
Since $V_4$ has exponent two, $2\mathrm T$ lies in the translation
kernel and $\rho_{\mathrm T}$ factors through
$\mathrm T/2\mathrm T\simeq\mathbb Z_2^3$. It can be written as
\stepcounter{equation}
\begin{equation}
 \rho_{\mathrm T}(\bm t)=
 \operatorname{diag}(e^{i\bm k_1\cdot\bm t},
 e^{i\bm k_2\cdot\bm t},e^{i\bm k_3\cdot\bm t}),\qquad\bm t\in L,
 \tag{24}\label{algmerge:kleintranslation}
\end{equation}
where the $\bm k_i$ are TRIMs,
\stepcounter{equation}
\begin{equation}
 \bm k_i=\tfrac12(n_{i1}\bm b_1+n_{i2}\bm b_2+n_{i3}\bm b_3),
 \qquad n_{ij}\in\{0,1\}.
 \tag{25}\label{algmerge:trim}
\end{equation}
The image is $V_4$ precisely when the three characters are distinct,
nonzero, and satisfy
\stepcounter{equation}
\begin{equation}
 \bm k_1+\bm k_2+\bm k_3=0\pmod{\widehat L}.
 \tag{26}\label{algmerge:sumtrim}
\end{equation}
Equivalently, $\ker\rho_{\mathrm T}=\langle2\mathrm T,\bm t_0\rangle$
for a nonzero $\bm t_0\in L/2L$. The three characters are exactly
its nonzero annihilators,
\stepcounter{equation}
\begin{equation}
 \bm k_i\cdot\bm t_0=0\pmod{2\pi},\qquad i=1,2,3.
 \tag{27}\label{algmerge:annihilators}
\end{equation}
Normality of the translation kernel further requires
\stepcounter{equation}
\begin{equation}
 R_p\bm t_0=\bm t_0\pmod{2L},\qquad p\in P.
 \tag{28}\label{algmerge:invariantkernel}
\end{equation}
In a fixed primitive basis, the seven possible nonzero vectors are
$\bm a_i$, $\bm a_i+\bm a_j$ with $i<j$, and
$\bm a_1+\bm a_2+\bm a_3$. They respectively describe one trivial
primitive spin translation, two identical primitive half-turns,
and three distinct primitive half-turns. For a fixed kernel all
orderings of the annihilator triple are related by spin-axis
permutations. The invariant kernels are retained individually until
spatial equivalence is imposed; basis-dependent distinctions among
these seven vectors are not additional physical invariants.

\subsubsection{Decomposition of \texorpdfstring{$\rho$}{rho}}

The normalizer of $V_4$ in $SO(3)$ is the chiral octahedral group,
\stepcounter{equation}
\begin{equation}
 N_{SO(3)}(V_4)=O,
 \tag{29}\label{algmerge:kleinnormalizer}
\end{equation}
and $O/V_4\simeq S_3$ permutes the three spin axes. Thus
\stepcounter{equation}
\begin{equation}
 \rho(g)=D(g)S(p),\qquad D(g)\in V_4,\qquad S(p)\in\widetilde S_3\subset SO(3).
 \tag{30}\label{algmerge:kleindecomposition}
\end{equation}
Here $\widetilde S_3$ is a specified lift of the abstract permutation
group. If $P_\sigma\bm e_i=\bm e_{\sigma(i)}$, choose
\stepcounter{equation}
\begin{equation}
 S_\sigma=\operatorname{sgn}(\sigma)P_\sigma,
 \qquad S_\sigma S_\tau=S_{\sigma\tau},\qquad\det S_\sigma=1.
 \tag{30a}\label{algmerge:properlift}
\end{equation}
An ordinary odd permutation matrix would have determinant $-1$
and is therefore not the required lift. Equation~\eqref{algmerge:properlift}
has the same conjugation action on diagonal $V_4$ and introduces
no extra factor system. With $\pi_p(d)=S(p)dS(p)^{-1}$, the
remaining factor obeys
\stepcounter{equation}
\begin{equation}
 D(g_1)\pi_{p_1}\bigl(D(g_2)\bigr)=D(g_1g_2).
 \tag{31}\label{algmerge:twistedhom}
\end{equation}
It is a $\pi$-twisted homomorphism; conjugation by an element of
$V_4$ changes it by a one-coboundary.

\subsubsection{Covariant condition between \texorpdfstring{$\rho_{\mathrm T}$}{rho T} and \texorpdfstring{$S$}{S}}

Equation~(6) reads
\stepcounter{equation}
\begin{equation}
 \pi_p\bigl(\rho_{\mathrm T}(\bm t)\bigr)=\rho_{\mathrm T}(R_p\bm t).
 \tag{32}\label{algmerge:kleincovariance}
\end{equation}
In Cartesian coordinates, it implies
\stepcounter{equation}
\begin{equation}
 \bm k_{\sigma_p^{-1}(i)}=R_p^{-1}\bm k_i\pmod{\widehat L},
 \qquad i=1,2,3,
 \tag{33}\label{algmerge:permutationinverse}
\end{equation}
or
\stepcounter{equation}
\begin{equation}
 R_p\bm k_i=\bm k_{\sigma_p(i)}\pmod{\widehat L}.
 \tag{34}\label{algmerge:permutation}
\end{equation}
In reciprocal primitive coordinates this action is $M_p^{-T}$.
For an invariant kernel, its annihilator triple is preserved and
these equations determine $\sigma_p$ uniquely. We use the lift
\eqref{algmerge:properlift} to obtain $S(p)$. Thus $S$ is not an
independent choice once the translation holonomy is fixed.

\subsubsection{Obstructed 1-cocycle equation of point-group sections}

For fixed $\rho_{\mathrm T}$ and $S$, Eq.~(7) becomes
\stepcounter{equation}
\begin{equation}
 D(s(p_1))\pi_{p_1}\bigl(D(s(p_2))\bigr)
 =\rho_{\mathrm T}(\bm t_{p_1,p_2})D(s(p_1p_2)).
 \tag{35}\label{algmerge:type2cocycle}
\end{equation}
Write
\stepcounter{equation}
\begin{equation}
 D(s(p))=\operatorname{diag}(e^{i\theta_1^s(p)},e^{i\theta_2^s(p)},
 e^{i\theta_3^s(p)}),\qquad\theta_i^s(p)\in\{0,\pi\}.
 \tag{36}\label{algmerge:binaryphases}
\end{equation}
The determinant condition is
\stepcounter{equation}
\begin{equation}
 \theta_1^s(p)+\theta_2^s(p)+\theta_3^s(p)=0\pmod{2\pi},
 \tag{37}\label{algmerge:detconstraint}
\end{equation}
and Eq.~\eqref{algmerge:type2cocycle} reduces to
\stepcounter{equation}
\begin{equation}
 \theta_i^s(p_1)+\theta_{\sigma_{p_1}^{-1}(i)}^s(p_2)
 =\bm k_i\cdot\bm t_{p_1,p_2}+\theta_i^s(p_1p_2)\pmod{2\pi}.
 \tag{38}\label{algmerge:binarycocycle}
\end{equation}
Put $d_{p,i}=\theta_i^s(p)/\pi$ and
$b_{p_1p_2,i}=\bm k_i\cdot\bm t_{p_1,p_2}/\pi$ modulo two.
Equations~\eqref{algmerge:detconstraint} and
\eqref{algmerge:binarycocycle} are the finite linear system
\stepcounter{equation}
\begin{equation}
 d_{p_1,i}+d_{p_2,\sigma_{p_1}^{-1}(i)}-d_{p_1p_2,i}
 =b_{p_1p_2,i},\qquad \sum_i d_{p,i}=0,
 \qquad\text{over }\mathbb F_2.
 \tag{38a}\label{algmerge:binarysystem}
\end{equation}
Gaussian elimination decides whether this system is consistent.
A zero row with nonzero right-hand side detects the obstruction.
Otherwise, one particular solution $d^{(0)}$ and a basis
$h_1,\ldots,h_r$ of the homogeneous kernel give every solution:
\stepcounter{equation}
\begin{equation}
 d=d^{(0)}+\sum_{j=1}^r\epsilon_j h_j,
 \qquad\epsilon_j\in\mathbb F_2.
 \tag{38b}\label{algmerge:binarysolutions}
\end{equation}
Conjugation by a fixed diagonal element with entries $(-1)^{a_i}$,
$\sum_i a_i=0$, changes
$d_{p,i}$ by $a_i+a_{\sigma_p^{-1}(i)}$. Quotienting by these
coboundaries gives the torsor for $H^{1,\pi}(P,V_4)$, and the
remaining axis permutations are included in the full spin quotient.
Equivalently, one can construct the matrices from all solutions and
compare their complete sets of translations and sections under a
single signed permutation. The latter procedure is used here and
removes both kinds of redundancy at once.

\subsubsection{Algorithm}

The Type 2 enumeration is therefore:
\begin{enumerate}
\renewcommand{\labelenumi}{(\roman{enumi})}
\item Determine every nonzero $\bm t_0\in L/2L$ satisfying
Eq.~\eqref{algmerge:invariantkernel}, and construct its three
annihilating characters and the translation matrices.
\item Determine $\sigma_p$ from Eq.~\eqref{algmerge:permutation}
and use the orientation-preserving lift~\eqref{algmerge:properlift}.
\item Solve Eq.~\eqref{algmerge:binarysystem} over $\mathbb F_2$,
construct $D(s(p))$ for every affine solution, and obtain
$\rho(t s(p))=\rho_{\mathrm T}(t)D(s(p))S(p)$.
\item Identify spin-conjugate representations by simultaneous
signed permutations of all their matrices.
\item Combine with $\chi$ and impose the spatial equivalences
on the complete representation.
\end{enumerate}
For the coplanar branches in Sec.~III, the same binary procedure
uses two characters and omits the determinant-one row. It retains
all three unordered character pairs for each Klein kernel and the
pair $(0,\bm k)$ for the single-reflection image, as specified there.

\subsection{Equivalence relations}
\label{algmerge:equivalences}

Three equivalences are imposed: a global change of spin basis, an
allowed change of spatial coordinates, and continuation of a general
representation branch. The first two compare representations at
specified parameters. The third assigns a single family label to
representations with continuously variable parameters.

\subsubsection{Spin-space conjugation}

The spin equivalence is $\varrho_1(g)=Q\varrho_2(g)Q^{-1}$ for all
$g\in G$, with one common $Q\in O(N)$. For $N=3$, $Q$ and $-Q$
give the same conjugation, so proper rotations suffice. In Type 0,
the quotient construction above already removes the spin-basis
choices. In Type 1, the continuous gauge~\eqref{algmerge:gauge}
and the simultaneous reversal $(\bm k,\theta)\mapsto(-\bm k,-\theta)$
give the residual spin equivalence after the axis has been fixed.
In Type 2, the joint translation eigenspaces fix three axes, and
the remaining transformations are signed permutations. In the
coplanar problem they are restricted to $O(2)$ within the spin
plane. All matrices must be transformed together; comparing each
matrix under a different change of basis is insufficient.

\subsubsection{Spatial automorphisms}
\label{algmerge:spatial}

We use the chirality-preserving affine coordinate changes that
normalize the parent space group. Express the spatial matrices and
section translations in a primitive basis. An allowed $\{A\mid\bm o\}$
satisfies
\stepcounter{equation}
\begin{equation}
 A\in SL(3,\mathbb Z),\qquad
 A M_p A^{-1}=M_{\sigma(p)},\qquad
 \bm u_p=A\bm\tau_p+(I-M_{\sigma(p)})\bm o-\bm\tau_{\sigma(p)}
 \in\mathbb Z^3.
 \tag{38c}\label{algmerge:affine}
\end{equation}
The allowed coordinate group can be infinite, but only its
generators are needed. We first obtain the integral normalizer of
$P$. For triclinic groups it is generated by elementary integral
shears and a reflection. For monoclinic groups the repeated plane
and the distinct axis give $GL(2,\mathbb Z)\times\{\pm1\}$;
one takes the stabilizer of the actual primitive lattice in these
coordinates. Its orbit is finite because the lattice lies between
$\mathbb Z^3$ and $D^{-1}\mathbb Z^3$ for a fixed integer $D$.
This retains the plane--axis coupling in centered lattices.
For three distinct invariant axes, enumerate their signed
permutations and retain the integral normalizers. For an
irreducible plane plus an axis, or an irreducible three-dimensional
point representation, the normalizer preserves
$Q=\sum_p M_p^T M_p$. Its integral isometries are obtained from
the finite column bounds
$|v_i|^2\le Q_{jj}(Q^{-1})_{ii}$ and the required column inner
products. Taking the determinant kernel gives the proper
normalizer; merely discarding improper generators would not
necessarily generate that kernel.

It remains to preserve the nonsymmorphic extension. Stack the
section translations into $\tau$ and the matrices $I-M_p$ into
$B_0$. If the rows of $C_0$ form the full integer annihilator of
$B_0$, then $C_0\tau$ modulo one specifies the section array
modulo origin shifts. A linear normalizer acts by
$\tau^A_{\sigma(p)}=A\bm\tau_p$. Because the translations have
a fixed common denominator, the orbit of this class is finite.
Traverse that orbit under all proper linear generators. If
$t_x$ takes the initial class to $x$ and a generator $s$ takes
$x$ to $y$, the transformations
\stepcounter{equation}
\begin{equation}
 h=t_y^{-1}s t_x,\qquad
 B_0\bm o=\tau-\tau^h\pmod1
 \tag{38d}\label{algmerge:stabilizer}
\end{equation}
give generators of its stabilizer and their affine lifts.
The same finite-orbit construction gives the lattice stabilizer
and determinant kernel used above. It is necessary to use the
whole orbit, since products of nonpreserving generators may
preserve the extension. Finally solve $B_0\bm o=0$ modulo one
and include all discrete origin-shift components. Continuous
directions with $B_0\bm v=0$ centralize $G$; integer origin shifts
give inner automorphisms and hence only spin conjugations.
Both act trivially after the spin quotient.

The resulting automorphism acts by
$\alpha(t_{\bm n}s(p))=t_{A\bm n+\bm u_p}s(\sigma(p))$ and
pullback $\varrho\mapsto\varrho\circ\alpha$. On the Type 1
variables of Eq.~\eqref{algmerge:joint} this is
\stepcounter{equation}
\begin{equation}
 \bm q'=A^T\bm q,\qquad
 \eta'_p=\eta_{\sigma(p)},\qquad
 x'_p=x_{\sigma(p)}+\bm q\cdot\bm u_p\pmod1.
 \tag{38e}\label{algmerge:pullback}
\end{equation}
The gauge is then restored, and $\chi$ is pulled back by the
same $\alpha$. We compare the resulting full affine tori,
not only their wave-vector projections. Two components are
identified when their complete images coincide after the
allowed spin normalization. Applying the generators and merging
the corresponding sets until no further mergers occur gives
the finite orbit catalogue. Taking separate spatial quotients
of $\rho$ and $\chi$ would generally give an incorrect result.

\subsubsection{Continuous families}
\label{algmerge:continuous}

The family relation is defined by continuation of the complete
general representation branch with its discrete constituent
data fixed. It is not determined merely by equality of $\eta$
or by connectedness of a wave-vector domain. For each
positive-dimensional component $C_a$ in
Eq.~\eqref{algmerge:jointcomponents}, remove the loci
$2\bm q=0$ modulo one from its general representatives and
assign the remaining domain a single family label. Its removed
points are recorded in Type 0 or the nonzero-TRIM strata.
The isolated non-TRIM solutions are discrete, not continuous
families.

To understand why disconnected pieces of the remaining domain
can belong to one family, consider the projection
$C_a\to K_a$ onto its wave-vector image. At fixed wave vector,
an infinitesimal phase variation obeys
$z_{p_1}+\eta_{p_1}z_{p_2}-z_{p_1p_2}=0$. Averaging over
$p_2$ yields
\stepcounter{equation}
\begin{equation}
 z_p=(1-\eta_p)\bar z,\qquad
 \bar z=|P|^{-1}\sum_{q\in P}z_q.
 \tag{38f}\label{algmerge:phasefiber}
\end{equation}
This is exactly the rotation gauge already removed. The
remaining fiber is finite, so $C_a\to K_a$ is a connected
finite covering. Motion in this covering retains the entire
section representation, including its phase sheets.

For example, the axial branch of $P2_1$ with $\eta=1$ obeys
$2x_s=q_y$ modulo one. Its complete solution is
$(q_y,x_s)=(2u,u)$ modulo one, with $u\in\mathbb T$.
The two phases at fixed $q_y$ are exchanged by winding once
around the reciprocal period. Deleting the zero and half-period
wave vectors splits the parameter circle into arcs, but does
not produce different continuation families. A special value
may be crossed using the same branch formula, including a
zero wave vector. It does not thereby become a general member,
and it does not authorize switching to an unrelated branch
when its limiting matrices happen to coincide.

\subsubsection{Equivalence with Song's representation classification}

With the same chirality-preserving coordinate convention and
the continuation convention just specified, the above quotient
is equivalent to Song's representation classification.\footnote{%
Z. Xiao, J. Zhao, Y. Li, R. Shindou, and Z.-D. Song,
\textit{Spin Space Groups: Full Classification and Applications},
Phys. Rev. X \textbf{14}, 031037 (2024), Sec. II D and Appendix B.}
This can be seen directly at the level of the generating
equivalences. Both descriptions classify the same real
orthogonal representations; decomposition into irreducible
constituents and restriction to translation holonomy are two
ways of specifying them. A change of spin basis is the same
simultaneous conjugation in either description, and an allowed
spatial transformation is the same pullback by
Eq.~\eqref{algmerge:affine}.

For a general continuous representation, the translation
characters occur as a conjugate pair. Its real two-dimensional
block is the realification of a complex character when
$\eta=1$, and has complexification induced from that character
on the index-two subgroup $\ker\eta$ when $\eta$ is nontrivial.
These are the continuous two-dimensional blocks used in Song's
classification. Conversely, choosing a real orthogonal basis
for either block gives Eq.~\eqref{algmerge:axialsection} and
the joint equations~\eqref{algmerge:joint}. In three dimensions
the remaining real character must be retained. More explicitly,
if $\varrho=\beta\oplus\nu$, where $\beta$ is the plane block,
then
\stepcounter{equation}
\begin{equation}
 \chi=(\det\beta)\nu,\qquad D=\chi\beta,\qquad
 \varrho=\chi\operatorname{diag}(D,\det D).
 \tag{38g}\label{algmerge:dictionary}
\end{equation}
Thus the wave vector of $\beta$ differs from that of $D$ by
the translation character of $\chi$. This conversion is
invertible and commutes with spatial pullback.

The continuous operations in Song's Appendix B are motion
along a fixed representation branch, exchange of its conjugate
wave vectors, and winding along independent reciprocal periods.
Each lifts to a path in a component $C_a$, possibly followed
by $(\bm q,x)\mapsto(-\bm q,-x)$, and is therefore included
here. Conversely, the monodromy of a connected covering
$C_a\to K_a$ acts transitively on its phase sheets: a path
between two points over the same wave vector projects to a
closed loop exchanging those sheets. Since
$\pi_1(K_a)\simeq\mathbb Z^{d_a}$, these exchanges are generated
by winding along a primitive set of reciprocal periods.
Together with motion along the branch and the conjugate-pair
exchange, they generate every identification within $C_a$.
Special values are used only to continue that same branch;
the special representations themselves remain in their
separate types. These arguments give both directions of the
continuous correspondence, rather than inferring it from
matching class counts.

Finally, pullback maps a representation branch to another
branch and commutes with its continuation. Taking the
transitive closure of spin conjugation, branch continuation,
and the common spatial generators therefore gives the same
representation classes in the two constructions. This
equivalence concerns the classification prescription; members
of one continuous class need not be conjugate at arbitrary
specified parameter values.

% The next equation of the unchanged Section III must remain (39).
\endgroup
\setcounter{equation}{38}
